\documentclass[a4paper,12pt]{article}
\usepackage[utf8]{inputenc}
\usepackage[T1]{fontenc}
\usepackage[french]{babel}
\usepackage{amsmath, amssymb, amsfonts}
\usepackage{array, longtable, booktabs}
\usepackage[table]{xcolor} % option table pour \rowcolors
\usepackage{geometry}
\geometry{left=1.8cm, right=1.8cm, top=2cm, bottom=2cm}
\usepackage{enumitem}
\usepackage{hyperref}
\hypersetup{colorlinks=true, linkcolor=blue, urlcolor=blue}

% Commandes pour les coches
\newcommand{\fait}{\ensuremath{\checkmark}}
\newcommand{\revoit}{\ensuremath{\circlearrowright}}
\newcommand{\casevide}{\ensuremath{\square}}

% Types de colonnes
\newcolumntype{C}[1]{>{\centering\arraybackslash}p{#1}}
\newcolumntype{L}[1]{>{\raggedright\arraybackslash}p{#1}}
\newcolumntype{R}[1]{>{\raggedleft\arraybackslash}p{#1}}

% Pour réduire les débordements dans les longtable
\setlength{\LTleft}{-0.5cm}
\setlength{\LTright}{-0.5cm}

\title{Programmation annuelle de Mathématiques -- 1\up{re} STI2D (spécialité PCM) \\
	Année 2026--2027 -- Zone C (3 h/semaine)}
\author{}
\date{Version du \today}

\begin{document}
	
	\maketitle
	
	\section*{Présentation générale}
	Ce document propose une organisation spiralaire de l’enseignement de la partie mathématiques de la spécialité « Physique-Chimie et Mathématiques » en première STI2D, conforme au programme officiel. L’année comporte 31 semaines de cours effectifs (hors vacances et jours fériés). Chaque notion est introduite une première fois, puis réinvestie et approfondie dans un second temps, en lien avec les autres thèmes et avec la physique-chimie. Cette approche favorise une acquisition progressive et durable des connaissances.
	
	\medskip
	\textbf{Place de Python :} La programmation est intégrée de façon transverse et évaluée dans \emph{tous} les sujets (formatives, contrôles, DM, PBL). Les compétences algorithmiques (boucles, tests, fonctions, simulation, calcul approché) sont développées en parallèle des mathématiques et réactivées régulièrement selon le même principe spiralé.
	
	\medskip
	\textbf{Utilisation de l’IA dans les DM :} Dans chaque devoir maison, l’élève est invité à utiliser un outil d’intelligence artificielle générative (ChatGPT, Mistral, Copilot, etc.) pour :
	\begin{itemize}
		\item obtenir des explications ou des exemples supplémentaires sur une notion,
		\item suggérer des pistes de résolution,
		\item débuguer ou améliorer un code Python,
		\item confronter différentes méthodes de résolution.
	\end{itemize}
	\emph{L’IA ne doit cependant pas rédiger la solution finale à la place de l’élève.} Le raisonnement personnel et la rédaction restent exigibles. Un \textbf{journal de bord} (questions posées, réponses obtenues, critiques et adaptations apportées) est à annexer à chaque DM. Cette pratique développe l’esprit critique et la capacité à évaluer la pertinence des réponses générées.
	
	\section{Tableau récapitulatif des contenus à enseigner}
	Le tableau ci-dessous reprend l’intégralité des attendus du programme de mathématiques, enrichi des éléments issus de la synthèse de trigonométrie. La colonne \textbf{« Traité »} indique si le contenu a été abordé en classe (coche \fait), et la colonne \textbf{« Reprise »} signale les notions qui feront l’objet d’un second passage (\revoit) pour consolider la compréhension.
	
	\rowcolors{2}{gray!8}{white}
	\begin{longtable}{|L{13.5cm}|C{1.5cm}|C{1.5cm}|}
		\hline
		\textbf{Thème / Contenu} & \textbf{Traité} & \textbf{Reprise} \\
		\hline
		\endhead
		\hline
		\multicolumn{3}{r}{Suite de la page suivante} \\
		\endfoot
		\hline
		\endlastfoot
		
		% --- Trigonométrie : approche par les transformations ---
		\multicolumn{3}{|L{13.5cm}|}{\textbf{Trigonométrie – Approche par les transformations}} \\
		Cercle trigonométrique : coordonnées $(x,y)$ d’un point, lien avec $\cos\theta$, $\sin\theta$ & \casevide & \casevide \\
		Transformations du plan : symétrie axiale (Ox, Oy), symétrie centrale, rotation d’angle $\pm\pi/2$ & \casevide & \casevide \\
		\multicolumn{3}{|L{13.5cm}|}{\textit{Lien avec les angles associés et les formules correspondantes}} \\
		Symétrie axiale par rapport à l’axe des ordonnées : $\cos(\pi-x)=-\cos x$, $\sin(\pi-x)=\sin x$ & \casevide & \casevide \\
		Symétrie centrale : $\cos(\pi+x)=-\cos x$, $\sin(\pi+x)=-\sin x$ & \casevide & \casevide \\
		Symétrie axiale par rapport à l’axe des abscisses : $\cos(-x)=\cos x$, $\sin(-x)=-\sin x$ & \casevide & \casevide \\
		Rotation de $+\pi/2$ : $\cos(\pi/2+x)=-\sin x$, $\sin(\pi/2+x)=\cos x$ & \casevide & \casevide \\
		Rotation de $-\pi/2$ : $\cos(\pi/2-x)=\sin x$, $\sin(\pi/2-x)=\cos x$ & \casevide & \casevide \\
		Valeurs remarquables : 0, $\pi/6$, $\pi/4$, $\pi/3$, $\pi/2$, $\pi$ (déduites de ces transformations) & \casevide & \casevide \\
		\hline
		
		% --- Fonctions trigonométriques ---
		\multicolumn{3}{|L{13.5cm}|}{\textbf{Fonctions trigonométriques}} \\
		Fonctions sinus et cosinus : périodicité, variations, parité & \casevide & \casevide \\
		Fonctions $A\cos(\omega t+\varphi)$ et $A\sin(\omega t+\varphi)$ : amplitude, période, phase & \casevide & \casevide \\
		Conversions degrés/radians, résolution d’équations $\cos x=a$, $\sin x=a$ & \casevide & \casevide \\
		\hline
		
		% --- Produit scalaire ---
		\multicolumn{3}{|L{13.5cm}|}{\textbf{Produit scalaire}} \\
		Définition géométrique, projection orthogonale, propriétés (bilinéarité, symétrie) & \casevide & \casevide \\
		Expression dans une base orthonormée, orthogonalité, calcul de longueurs/angles & \casevide & \casevide \\
		Théorème d’Al-Kashi, égalité du parallélogramme & \casevide & \casevide \\
		\hline
		
		% --- Nombres complexes ---
		\multicolumn{3}{|L{13.5cm}|}{\textbf{Nombres complexes}} \\
		Forme algébrique, conjugué, module, représentation dans le plan, affixe d’un point/vecteur & \casevide & \casevide \\
		Somme, produit, quotient, conjugué d’une somme/produit/quotient, module d’un produit/quotient & \casevide & \casevide \\
		Argument, forme trigonométrique (sans exponentielle) & \casevide & \casevide \\
		Passage algébrique $\leftrightarrow$ trigonométrique, interprétation géométrique & \casevide & \casevide \\
		\hline
		
		% --- Dérivées ---
		\multicolumn{3}{|L{13.5cm}|}{\textbf{Dérivées}} \\
		Taux de variation, nombre dérivé, notations, approximation affine & \casevide & \casevide \\
		Dérivée d’une somme, d’un produit, de l’inverse, d’un quotient & \casevide & \casevide \\
		Dérivée de $x \mapsto x^n$, de $x \mapsto 1/x$, dérivée des polynômes & \casevide & \casevide \\
		Dérivée de $\cos$ et $\sin$, de $f(ax+b)$, de $A\cos(\omega t+\varphi)$ et $A\sin(\omega t+\varphi)$ & \casevide & \casevide \\
		Étude du sens de variation d’une fonction & \casevide & \casevide \\
		\hline
		
		% --- Primitives ---
		\multicolumn{3}{|L{13.5cm}|}{\textbf{Primitives}} \\
		Définition d’une primitive, deux primitives diffèrent d’une constante & \casevide & \casevide \\
		Primitives d’un polynôme, de $A\cos(\omega t+\varphi)$ et $A\sin(\omega t+\varphi)$ & \casevide & \casevide \\
		Méthode d’Euler pour l’approximation d’une primitive & \casevide & \casevide \\
		\hline
	\end{longtable}
	
	\section{Calendrier et programmation harmonieuse}
	
	\subsection{Calendrier retenu (zone C, jours ouvrés)}
	Les vacances et jours fériés ont été pris en compte. Les semaines de cours sont les suivantes :
	
	\begin{center}
		\begin{tabular}{|C{3.5cm}|C{5cm}|C{2.5cm}|}
			\hline
			\textbf{Période} & \textbf{Dates} & \textbf{Nb de semaines} \\
			\hline
			Rentrée – Toussaint & 1er sept. – 16 oct. 2026 & 7 \\
			Toussaint – Noël & 2 nov. – 18 déc. 2026 & 7 \\
			Noël – Hiver & 4 janv. – 5 fév. 2027 & 5 \\
			Hiver – Printemps & 22 fév. – 2 avr. 2027 & 6 \\
			Printemps – Fin mai & 19 avr. – 31 mai 2027 & 6 \\
			\hline
			\textbf{Total} & & \textbf{31 semaines} \\
			\hline
		\end{tabular}
	\end{center}
	
	\textit{Jours fériés (tombant un jour de cours) exclus :} 1er janvier 2027 (vendredi), 29 mars 2027 (lundi de Pâques), 6 mai 2027 (Ascension, jeudi), 17 mai 2027 (Pentecôte, lundi). Les 1er et 8 mai 2027 tombent un samedi.
	
	\subsection{Programmation par trimestre (approche spiralée)}
	Les tableaux suivants précisent pour chaque semaine :
	\begin{itemize}
		\item les contenus nouveaux,
		\item les \textbf{points ciblés dans l’évaluation formative} (5 à 10 min en début de séance) incluant systématiquement une question Python,
		\item des remarques utiles.
	\end{itemize}
	Les évaluations formatives mêlent mathématiques et programmation, avec réactivation des notions antérieures.
	
	% --- 1er trimestre ---
	\newpage
	\subsubsection{1\textsuperscript{er} trimestre (septembre – novembre 2026) – 14 semaines}
	
	\rowcolors{2}{gray!8}{white}
	\begin{longtable}{|C{0.7cm}|L{3.8cm}|L{5.5cm}|L{2.3cm}|}
		\hline
		\textbf{Sem.} & \textbf{Contenus nouveaux} & \textbf{Évaluation formative -- points évalués (dont Python)} & \textbf{Remarques} \\
		\hline
		\endhead
		\hline
		\multicolumn{4}{r}{Suite de la page suivante} \\
		\endfoot
		\hline
		\endlastfoot
		
		1-2 & \textbf{Trigonométrie – approche par les transformations du plan.} 
		Cercle trigonométrique, coordonnées d’un point, lien avec $\cos\theta$ et $\sin\theta$.  
		Réactivation : symétries axiales (Ox, Oy), symétrie centrale, rotations de $\pm\pi/2$ (vues en 2nde).  
		Application aux formules d’angles associés. & Calculs de coordonnées de points après transformation.  
		\textbf{Python} : tracer un cercle trigonométrique, placer un point et ses transformés par symétries/rotations. & Point de départ : les élèves connaissent déjà les transformations du plan. On en déduit les formules trigonométriques. \\
		\hline
		3-4 & Établissement des formules d’angles associés :  
		$\cos(-x)$, $\cos(\pi-x)$, $\cos(\pi+x)$, $\cos(\pi/2+x)$, $\cos(\pi/2-x)$ ; idem pour $\sin$.  
		Détermination des valeurs remarquables (0, $\pi/6$, $\pi/4$, $\pi/3$, $\pi/2$, $\pi$). & \textbf{Spirale :} transformations (S1-2).  
		\textbf{Python} : créer un programme qui, étant donné un angle, affiche son image par chaque transformation. & Les formules sont justifiées géométriquement. \\
		\hline
		5-6 & Fonctions $A\cos(\omega t+\varphi)$ et $A\sin(\omega t+\varphi)$ : amplitude, période, phase. & \textbf{Spirale :} valeurs remarquables (S3-4).  
		\textbf{Python} : tracé de courbes sinusoïdales avec \texttt{matplotlib}. & Lien avec les signaux électriques (physique). \\
		\hline
		7-8 & Résolution d’équations $\cos x=a$, $\sin x=a$. Périodicité et symétries du cercle. & \textbf{Spirale :} fonctions sinusoïdales (S5-6).  
		\textbf{Python} : résolution numérique par balayage. & Utilisation du cercle trigonométrique. \\
		\hline
		9-10 & Produit scalaire : définition géométrique, projection orthogonale, propriétés (bilinéarité, symétrie). & \textbf{Spirale :} trigonométrie (S1-8).  
		\textbf{Python} : calcul de produit scalaire avec listes. & Lien avec le travail d’une force (physique). \\
		\hline
		11-12 & Expression du produit scalaire en base orthonormée, orthogonalité, calcul de longueurs/angles. Théorème d’Al-Kashi. & \textbf{Spirale :} produit scalaire (S9-10).  
		\textbf{Python} : calcul de norme et d’angle. & Application à la mécanique. \\
		\hline
		13-14 & Nombres complexes : forme algébrique, conjugué, module, représentation. & \textbf{Bilan T1 :} trigo, produit scalaire, début complexes.  
		\textbf{Python} : module et conjugué avec \texttt{cmath}. & Introduction aux complexes. \\
		\hline
		\multicolumn{4}{l}{} \\
		\hline
		\multicolumn{4}{l}{\textbf{Contrôles continus :} semaines 5, 10, 14 (avec exercice Python). \textbf{Formatives :} chaque semaine.} \\
		\hline
	\end{longtable}
	
	% --- 2e trimestre ---
	\newpage
	\subsubsection{2\textsuperscript{e} trimestre (novembre 2026 – février 2027) – 12 semaines}
	
	\rowcolors{2}{gray!8}{white}
	\begin{longtable}{|C{0.7cm}|L{3.8cm}|L{5.5cm}|L{2.3cm}|}
		\hline
		\textbf{Sem.} & \textbf{Contenus nouveaux} & \textbf{Évaluation formative -- points évalués (dont Python)} & \textbf{Remarques} \\
		\hline
		\endhead
		\hline
		\multicolumn{4}{r}{Suite de la page suivante} \\
		\endfoot
		\hline
		\endlastfoot
		
		15-16 & Nombres complexes : opérations (somme, produit, quotient), conjugué d’une somme/produit/quotient, module d’un produit/quotient. & \textbf{Spirale :} forme algébrique (S13-14).  
		\textbf{Python} : fonctions pour calculer conjugué, module. & Calculs algébriques. \\
		\hline
		17-18 & Argument, forme trigonométrique d’un complexe, passage algébrique $\leftrightarrow$ trigonométrique. & \textbf{Spirale :} opérations sur les complexes (S15-16).  
		\textbf{Python} : conversion avec \texttt{cmath}. & Lien avec la trigonométrie. \\
		\hline
		19-20 & Interprétation géométrique des complexes (affixe, distance, angle). & \textbf{Spirale :} forme trigonométrique (S17-18).  
		\textbf{Python} : visualisation dans le plan complexe. & \textbf{PBL 1 (2h)} : Étude de circuits électriques en régime sinusoïdal (utilisation des complexes). \\
		\hline
		21-22 & Dérivées : taux de variation, nombre dérivé, notations, approximation affine. & \textbf{Spirale :} fonctions trigonométriques (S1-4).  
		\textbf{Python} : calcul de taux de variation. & Lien avec la physique (vitesse). \\
		\hline
		23-24 & Dérivée d’une somme, d’un produit, de l’inverse, d’un quotient ; dérivée de $x^n$, de $1/x$. & \textbf{Spirale :} approximation affine (S21-22).  
		\textbf{Python} : fonction qui calcule la dérivée d’un polynôme. & Calcul formel. \\
		\hline
		25-26 & Dérivée des fonctions $\cos$ et $\sin$, dérivée de $f(ax+b)$, de $A\cos(\omega t+\varphi)$ et $A\sin(\omega t+\varphi)$. & \textbf{Spirale :} règles de dérivation (S23-24).  
		\textbf{Python} : dérivation de fonctions sinusoïdales. & \textbf{PBL 2 (2h)} : Modélisation d’un mouvement oscillatoire (ressort, pendule). \\
		\hline
		\multicolumn{4}{l}{} \\
		\hline
		\multicolumn{4}{l}{\textbf{Contrôles continus :} semaines 18, 22, 26 (avec exercice Python). \textbf{Formatives :} chaque semaine.} \\
		\hline
	\end{longtable}
	
	% --- 3e trimestre ---
	\newpage
	\subsubsection{3\textsuperscript{e} trimestre (février – mai 2027) – 12 semaines}
	
	\rowcolors{2}{gray!8}{white}
	\begin{longtable}{|C{0.7cm}|L{3.8cm}|L{5.5cm}|L{2.3cm}|}
		\hline
		\textbf{Sem.} & \textbf{Contenus nouveaux} & \textbf{Évaluation formative -- points évalués (dont Python)} & \textbf{Remarques} \\
		\hline
		\endhead
		\hline
		\multicolumn{4}{r}{Suite de la page suivante} \\
		\endfoot
		\hline
		\endlastfoot
		
		27-28 & Primitives : définition, deux primitives diffèrent d’une constante, primitives d’un polynôme. & \textbf{Spirale :} dérivation (S21-26).  
		\textbf{Python} : calcul de primitives de polynômes. & Lien avec l’intégration en physique (travail, énergie). \\
		\hline
		29-30 & Primitives de $t\mapsto A\cos(\omega t+\varphi)$ et $A\sin(\omega t+\varphi)$. & \textbf{Spirale :} primitives de polynômes (S27-28).  
		\textbf{Python} : vérification par dérivation. & Applications en électricité (puissance). \\
		\hline
		31-32 & Méthode d’Euler pour l’approximation d’une primitive (problème de Cauchy). & \textbf{Spirale :} dérivation, primitives (S21-30).  
		\textbf{Python} : implémentation de la méthode d’Euler. & \textbf{PBL 3 (2h)} : Résolution approchée d’une équation différentielle (ex. décroissance radioactive). \\
		\hline
		33-34 & Synthèse sur les nombres complexes, la trigonométrie et les dérivées/primitives. & \textbf{Bilan complexe + analyse}.  
		\textbf{Python} : projet de modélisation d’un phénomène oscillant. & Interdisciplinarité. \\
		\hline
		35-36 & Synthèse produit scalaire et géométrie. & \textbf{Bilan géométrie}.  
		\textbf{Python} : calculs de distances et d’angles. & Préparation au contrôle. \\
		\hline
		37-38 & Révisions et évaluations finales. & \textbf{Bilan général :} tous les thèmes + Python. & Préparation épreuves. \\
		\hline
		\multicolumn{4}{l}{} \\
		\hline
		\multicolumn{4}{l}{\textbf{Contrôles continus :} semaines 30, 34, 38 (avec exercice Python). \textbf{Formatives :} chaque semaine.} \\
		\hline
	\end{longtable}
	
	\section{Évaluations}
	
	\subsection{Évaluations formatives (hebdomadaires)}
	Chaque semaine, une courte activité (5–10 min) mêle mathématiques et Python. Les points évalués sont détaillés dans les tableaux ci-dessus. Ces évaluations ne sont pas notées, mais elles permettent une auto-évaluation et un suivi des compétences algorithmiques.
	
	\subsection{Évaluations sommatives (contrôles continus)}
	Au moins trois contrôles par trimestre, d’une durée de 1 h chacun.  
	Chaque contrôle comporte :
	\begin{itemize}
		\item une partie sans calculatrice (calculs, raisonnement),
		\item une partie avec calculatrice / Python (exercice de programmation : écriture de fonction, complétion de code, simulation, représentation graphique, etc.).
	\end{itemize}
	
	\begin{center}
		\begin{tabular}{|C{3.5cm}|C{3.5cm}|C{3.5cm}|C{3.5cm}|}
			\hline
			\textbf{Trimestre} & \textbf{Contrôle 1} & \textbf{Contrôle 2} & \textbf{Contrôle 3} \\
			\hline
			1\textsuperscript{er} (sem. 5, 10, 14) & Trigonométrie (angles associés, symétries, équations) \newline + Python (visualisations) & Produit scalaire (projection, Al-Kashi) \newline + Python (calculs) & Nombres complexes (algèbre, module) \newline + Python (conjugué, module) \\
			\hline
			2\textsuperscript{e} (sem. 18, 22, 26) & Complexes (forme trigonométrique, géométrie) \newline + Python (conversions) & Dérivées (règles, composées) \newline + Python (dérivée formelle) & Dérivées des fonctions trigonométriques \newline + Python (tracés, tangentes) \\
			\hline
			3\textsuperscript{e} (sem. 30, 34, 38) & Primitives (polynômes, sin/cos) \newline + Python (primitives) & Méthode d’Euler \newline + Python (implémentation) & Synthèse (complexes, dérivées, primitives) \newline + Python (projet complet) \\
			\hline
		\end{tabular}
	\end{center}
	
	\subsection{Problèmes basés sur l’apprentissage (PBL) – 2 h chacun}
	Trois PBL sont répartis sur l’année, en séances de 2 h (travail en groupes, compte rendu écrit et code Python rendu) :
	
	\begin{enumerate}
		\item \textbf{PBL 1 (semaine 20)} : \emph{Étude de circuits électriques en régime sinusoïdal}. Utilisation des nombres complexes pour représenter les impédances, calculer des courants et tensions. \textbf{Python} : simulation et tracé des diagrammes de Fresnel.
		\item \textbf{PBL 2 (semaine 26)} : \emph{Modélisation d’un mouvement oscillatoire} (ressort, pendule). Utiliser les fonctions $A\cos(\omega t+\varphi)$ et les dérivées pour étudier la position, vitesse, accélération. \textbf{Python} : tracé des courbes et analyse énergétique.
		\item \textbf{PBL 3 (semaine 32)} : \emph{Résolution approchée d’une équation différentielle} (ex. décroissance radioactive, refroidissement de Newton). Appliquer la méthode d’Euler pour approcher la solution et comparer avec la solution exacte. \textbf{Python} : implémentation d’Euler et visualisation.
	\end{enumerate}
	
	\subsection{PBL dédié à la trigonométrie – « Cercle trigonométrique et transformations interactives » (semaine 4)}
	Ce projet spécifique est proposé en semaine 4 pour approfondir la maîtrise des liens entre transformations du plan, coordonnées et angles associés :
	
	\begin{itemize}
		\item \textbf{Objectif :} Construire un cercle trigonométrique interactif avec Python (bibliothèque \texttt{matplotlib}) permettant de visualiser :
		\begin{itemize}
			\item les coordonnées $(x,y)$ d’un point d’angle $\theta$,
			\item les images de ce point par les transformations suivantes : symétrie axiale (Ox, Oy), symétrie centrale, rotations de $+\pi/2$ et $-\pi/2$,
			\item les formules d’angles associés correspondantes.
		\end{itemize}
		\item \textbf{Démarche :} Les élèves écrivent un script qui :
		\begin{enumerate}
			\item place le point $M(\cos\theta,\sin\theta)$ sur le cercle,
			\item affiche les coordonnées des points transformés,
			\item compare les coordonnées obtenues avec les formules théoriques,
			\item propose un curseur pour modifier $\theta$ en temps réel.
		\end{enumerate}
		\item \textbf{Évaluation :} Code rendu, accompagné d’une brève explication des choix réalisés.
	\end{itemize}
	
	\section{Devoirs maison (DM) -- progression, Python et utilisation de l’IA}
	
	Afin de favoriser une acquisition à long terme, un devoir maison est distribué \textbf{toutes les deux semaines} (18 DM sur l’année).  
	Chaque DM porte sur l’ensemble du programme traité depuis le 1er septembre, et comporte :
	\begin{itemize}
		\item des exercices de calcul et de raisonnement,
		\item un exercice de programmation Python,
		\item une mission spécifique d’utilisation de l’IA (consultation, débuggage, confrontation de méthodes ou critique de réponse).
	\end{itemize}
	
	L’élève doit joindre à sa copie un \textbf{journal de bord} retraçant ses échanges avec l’IA et les décisions prises à la suite de ces échanges.
	
	\begin{center}
		\small
		\rowcolors{2}{gray!8}{white}
		\begin{longtable}{|C{0.7cm}|C{0.9cm}|L{6.0cm}|L{6.0cm}|}
			\hline
			\textbf{DM} & \textbf{Sem.} & \textbf{Thèmes évalués (mathématiques + Python)} & \textbf{Mission IA spécifique} \\
			\hline
			\endhead
			\hline
			\multicolumn{4}{r}{Suite à la page suivante} \\
			\endfoot
			\hline
			\endlastfoot
			
			1 & 3 & Transformations du plan et angles associés : symétries, rotations, formules. Python : visualisation. & Demander à l’IA d’expliquer comment la symétrie axiale par rapport à Oy se traduit sur les coordonnées. \\
			\hline
			2 & 5 & Fonctions $A\cos(\omega t+\varphi)$, $A\sin(\omega t+\varphi)$. Python : tracé. & Utiliser l’IA pour générer une sinusoïde et ajuster ses paramètres. \\
			\hline
			3 & 7 & Résolution d’équations $\cos x=a$, $\sin x=a$, périodicité. Python : balayage. & Demander à l’IA de résoudre une équation et de justifier la méthode. \\
			\hline
			4 & 9 & Produit scalaire : définition, projection. Python : calcul de produit scalaire. & Utiliser l’IA pour interpréter le produit scalaire en termes de travail d’une force. \\
			\hline
			5 & 11 & Expression en base orthonormée, orthogonalité. Théorème d’Al-Kashi. Python : calcul de normes. & Demander à l’IA de démontrer l’orthogonalité de deux vecteurs donnés. \\
			\hline
			6 & 13 & Nombres complexes : forme algébrique, conjugué, module. Python : calculs. & Demander à l’IA d’expliquer l’utilité des complexes en électricité. \\
			\hline
			7 & 16 & Opérations sur les complexes, forme trigonométrique. Python : conversion. & Utiliser l’IA pour vérifier une conversion et discuter des erreurs possibles. \\
			\hline
			8 & 18 & Interprétation géométrique des complexes. Python : tracé dans le plan. & Demander à l’IA de proposer un problème de géométrie avec les complexes. \\
			\hline
			9 & 20 & Dérivées : taux de variation, approximation affine. Python : calcul de taux. & Utiliser l’IA pour expliquer l’approximation affine avec un exemple concret (vitesse). \\
			\hline
			10 & 22 & Règles de dérivation (somme, produit, inverse, quotient). Python : dérivée formelle. & Demander à l’IA de démontrer la dérivée d’un produit. \\
			\hline
			11 & 24 & Dérivée de $\cos$, $\sin$, de $f(ax+b)$. Python : dérivée de fonctions trigo. & Utiliser l’IA pour vérifier les dérivées et les appliquer à un mouvement harmonique. \\
			\hline
			12 & 26 & Primitives de polynômes. Python : calcul de primitives. & Demander à l’IA de justifier la constante d’intégration. \\
			\hline
			13 & 28 & Primitives de $A\cos(\omega t+\varphi)$, $A\sin(\omega t+\varphi)$. Python : vérification par dérivation. & Utiliser l’IA pour relier primitive et énergie en physique. \\
			\hline
			14 & 30 & Méthode d’Euler. Python : implémentation. & Demander à l’IA de comparer la méthode d’Euler avec d’autres méthodes d’approximation. \\
			\hline
			15 & 32 & Synthèse complexes et trigonométrie. Python : projet de modélisation. & Utiliser l’IA pour critiquer un code Python et proposer des améliorations. \\
			\hline
			16 & 34 & Synthèse analyse (dérivées, primitives). Python : projet complet. & Demander à l’IA un retour sur la structure du code et proposer des extensions. \\
			\hline
			17 & 36 & Synthèse générale. Python : projet complet. & Demander à l’IA de suggérer des applications en physique-chimie. \\
			\hline
			18 & 38 & Bilan général. Python : projet complet. & Demander à l’IA de synthétiser les liens entre les différents chapitres. \\
			\hline
		\end{longtable}
	\end{center}
	
	\textbf{Remarques :}
	\begin{itemize}
		\item Les DM sont à rendre une semaine après la distribution.
		\item La mission IA doit être clairement identifiée dans le journal de bord (copie d’écran ou retranscription des questions/réponses, suivie de l’analyse personnelle de l’élève).
		\item La correction est effectuée en classe, avec une attention particulière aux erreurs fréquentes, aux méthodes de résolution et aux bonnes pratiques d’utilisation de l’IA.
	\end{itemize}
	
	\section{Conseils pour la mise en œuvre}
	\begin{itemize}
		\item \textbf{Rituels} : consacrer 5 à 10 minutes par séance à des exercices de calcul (mental, littéral) et à de courts extraits de code Python (lecture, complétion, débogage).
		\item \textbf{Travail en îlots} : favoriser les échanges entre élèves lors des phases de recherche et des PBL.
		\item \textbf{Numérique} : utiliser Python régulièrement (au moins une séance sur deux). Insister sur la programmation modulaire (fonctions) et la documentation du code.
		\item \textbf{IA éthique} : former les élèves à formuler des requêtes précises et à évaluer de manière critique les réponses de l’IA (vérification des calculs, cohérence, limites). Insister sur le fait que l’IA est un \emph{outil d’aide}, pas un substitut au raisonnement personnel.
		\item \textbf{Différenciation} : proposer des exercices de difficulté progressive et des approfondissements pour les élèves plus à l’aise (notamment en programmation et en interaction avec l’IA).
		\item \textbf{Trace écrite} : veiller à ce que chaque élève dispose d’un cahier de cours structuré (définitions, propriétés, exemples) et d’un cahier de programmation regroupant les fonctions et scripts clés.
	\end{itemize}
	
	\section*{Conclusion}
	Cette programmation respecte les 31 semaines de cours disponibles et couvre l’intégralité du programme de mathématiques de la spécialité PCM en 1\up{re} STI2D. L’approche choisie pour la trigonométrie – partir des transformations du plan pour faire émerger les formules d’angles associés – permet de donner du sens aux relations trigonométriques et de les ancrer dans des connaissances déjà maîtrisées. L’intégration systématique de Python et de l’utilisation critique de l’IA dans les apprentissages et les évaluations garantit que les notions mathématiques, la programmation et le jugement critique sont solidement maîtrisés en fin d’année. Les évaluations variées (formatives, sommatives, PBL, DM) permettent de suivre la progression de chaque élève et de développer ses compétences de modélisation, de raisonnement, de communication, de codage et de collaboration avec les outils d’intelligence artificielle.
	
\end{document}